\section{Results}
\label{sec:results}

Numbers come first in each subsection and interpretation is separated from
them. Every macro-F1 is reported against its own pair's chance baseline, and
because that baseline is pair-dependent, nothing is averaged across pairs.
Unless stated otherwise, intervals are 95\% paired cluster bootstrap intervals
over target-test speakers and seeds.

Sections~\ref{sec:results-translation} and~\ref{sec:results-boundary} test
Proposition 1 and delimit it. Sections~\ref{sec:results-selection}
to~\ref{sec:results-aggregation} place that result inside the frozen pipeline
comparison it was found in. Robustness controls and diagnostics whose
interpretation we deliberately restrict are summarised in
Section~\ref{sec:results-robustness} and reported in full in the supplementary
material. Table values and figure data are generated from the retained result
files, and both the original outcome report and the retrospective translation
report are included in the automated manuscript number-trace check.

\subsection{Exact validation equality in the frozen translation pairs}
\label{sec:results-translation}

Table~\ref{tab:translation-audit} tests the stored consequences of
Proposition 1. The audit joins \texttt{none} and \texttt{mean\_shift} by
direction, seed, backbone, layer, classifier, split identifiers and sizes,
and label-map, feature and search hashes. It requires one successful row per
arm and rejects duplicates or unmatched rows. No target score participates
in the join or the inclusion rule. The configuration and input-ledger hashes,
all matched run identifiers and unrounded values are released with the audit.

\begin{table*}[!t]
\centering\small
\caption{Retrospective translation audit on frozen runs. Each cell is one matched \texttt{none}/\texttt{mean\_shift} pair sharing direction, seed, backbone, layer setting, classifier and every split, label-map, feature and search identifier. Equal validation and equal parameters count exact equality of the stored winning source-validation score and the selected classifier parameters. $\Delta$ is mean target macro-F1 for mean shift minus none; Positive counts cells with a target improvement. Chance is the uniform-random baseline. The upper panel is the exact prediction of Proposition 1. The lower panel reports families whose fitting rules fall outside its assumptions, which the proposition does not cover and which are therefore scope evidence rather than counterexamples. Cells share seeds; no independence-based test is reported and no general claim about mean shift follows from the $\Delta$ column.}
\label{tab:translation-audit}
\begin{tabular}{llrrrrrr}
\toprule
Direction & Head & Pairs & Equal val. & Equal params & $\Delta$ & Positive & Chance \\
\midrule
\multicolumn{8}{l}{\textbf{Stationary kernel, Gram-determined fit: Proposition 1 applies}} \\
CREMA-D $\to$ RAVDESS & RBF SVM & 30 & 30 & 30 & +0.1969 & 30 & 0.1656 \\
RAVDESS $\to$ CREMA-D & RBF SVM & 30 & 30 & 30 & +0.1342 & 30 & 0.1665 \\
\midrule
\multicolumn{8}{l}{\textbf{Penalised intercept or iterative solver: outside those assumptions}} \\
CREMA-D $\to$ RAVDESS & Logistic & 30 & 0 & 22 & +0.1881 & 30 & 0.1656 \\
RAVDESS $\to$ CREMA-D & Logistic & 30 & 3 & 12 & +0.1442 & 30 & 0.1665 \\
CREMA-D $\to$ RAVDESS & Linear SVM & 30 & 0 & 19 & +0.1776 & 30 & 0.1656 \\
RAVDESS $\to$ CREMA-D & Linear SVM & 30 & 6 & 15 & +0.1456 & 30 & 0.1665 \\
CREMA-D $\to$ RAVDESS & MLP & 45 & 0 & 24 & +0.1701 & 45 & 0.1656 \\
RAVDESS $\to$ CREMA-D & MLP & 45 & 0 & 12 & +0.1287 & 45 & 0.1665 \\
\bottomrule
\end{tabular}
\\[2pt]
\begin{minipage}{\linewidth}\footnotesize Filter: \texttt{grid-freeze-v3}, \texttt{blending=none}; matched \texttt{none}/\texttt{mean\_shift} rows only. Every family spans three backbones and five seeds; the MLP additionally has the learned \texttt{weighted} layer aggregation, which the closed-form heads cannot use, hence its larger cell count. The reduced-seed transformer arm has no balanced matched cells and is excluded by construction. Generated by \texttt{tools/audit\_translation.py} from the configuration \texttt{configs/audit\_translation.yaml}. The JSON report retains input/configuration hashes and both run IDs for every pair. Means are descriptive, not independent-cell inference.\end{minipage}
\end{table*}


For the RBF SVM, all 30 pairs in each direction have exactly equal stored
winning validation scores and selected hyperparameters. This is the predicted
equality, observed without exception in the arm the proposition covers. Mean
shift improves target macro-F1 in every pair, with directional mean differences
of $+0.1342$ and $+0.1969$. Source-side validation therefore assigned identical
evidence to two conditions whose target predictions differ, which is the
failure the proposition describes.

The ledger stores the winning validation score for each run, not the
predictions of every search candidate. The audit is accordingly a check of the
recorded winning scores and selected parameters, not an empirical verification
of the entire search surface; the statement about every candidate remains
analytic. The analysis is retrospective and uses existing scores only, and
cells share seeds and speakers, so the table reports descriptive paired means
rather than a test treating each cell as an independent replication.

Target macro-F1 favours mean shift in every matched pair of every family. That
is a descriptive property of these frozen cells against an unaligned baseline
in one corpus pair. It is not evidence that mean matching improves cross-corpus
speech emotion recognition in general, and no such claim is made here.

\subsection{What the boundary classifiers show}
\label{sec:results-boundary}

The audit covers every classifier family in the frozen grid with complete,
balanced matched cells. Only the RBF SVM satisfies Proposition 1's fitting
assumptions, and Table~\ref{tab:translation-audit} is drawn in two panels on
that basis. The lower panel reports the three families the assumptions exclude,
identified in advance in Section~\ref{sec:translation-theory}: a penalised
synthetic intercept in the linear SVM, a finite iterative solver in logistic
regression, and initialisation, weight decay and early stopping in the MLP.

They behave as those assumptions predict. Logistic regression has exactly equal
stored validation scores in 3 forward pairs and none in reverse, with selected
hyperparameters agreeing in 12 forward and 22 reverse pairs. The linear SVM has
6 forward and no reverse exact equalities, with 15 and 19 parameter agreements.
The MLP, which additionally searches the learned \texttt{weighted} layer
aggregation and therefore has 45 pairs per direction, has no exact validation
equality in either direction and parameter agreement in 12 forward and 24
reverse pairs.

None of this contradicts Proposition 1, whose statement does not extend to a
penalised intercept or to iterative neural training. These rows are evidence
about the boundary of the proposition, not tests of it. The pattern separates
an exact identity between objectives from the numerical behaviour of the
solvers that approximate them, and it is reported rather than omitted so that
the scope of the proposition rests on the full set of eligible frozen cells.
The reduced-seed transformer arm has no balanced matched cells and is excluded
by construction.

\subsection{What the criterion costs a practitioner}
\label{sec:results-selection}

The invariance in Section~\ref{sec:results-translation} is exact but local: it
concerns two conditions the criterion cannot separate. Table~\ref{tab:headline}
gives the consequence for the whole protocol, comparing the configuration
selected on source-side validation alone against the best configuration present
in the grid. Its intervals are the one exception to the convention stated above:
because the selected configuration differs across seeds, this table reports a
$t$-interval over the five seeds rather than a paired cluster bootstrap.
Section~\ref{sec:metrics} states what that interval does and does not include.

\begin{table}[tb]
  \centering
  \caption{Target macro-F1 of the configuration selected on source validation, against the best configuration present in the grid. The oracle column is an upper bound no protocol can reach and is not a result.}
  \label{tab:headline}
  \begin{tabular}{lrrrrr}
    \toprule
    pair & validated & oracle & gap & chance & majority \\
    \midrule
    RAVDESS $\rightarrow$ CREMA-D & 0.3156 [0.1880, 0.4432] & 0.4555 [0.4326, 0.4785] & 0.1399 [0.0231, 0.2568] & 0.1665 & 0.0425 \\
    CREMA-D $\rightarrow$ RAVDESS & 0.4994 [0.4592, 0.5395] & 0.5680 [0.5148, 0.6212] & 0.0687 [0.0090, 0.1283] & 0.1656 & 0.0444 \\
    \bottomrule
  \end{tabular}
  \\[2pt]
  \begin{minipage}{\linewidth}\footnotesize Filter: \texttt{freeze\_tag=grid-freeze-v3}, \texttt{blending=none}, 4986 runs. Mean over 5 seeds with a 95\% $t$-interval. Both floors are analytic from the realised target-test priors.\end{minipage}
\end{table}


Validated macro-F1 is $0.3156$ $[0.1880, 0.4432]$ for
RAVDESS $\rightarrow$ CREMA-D against a chance baseline of $0.1665$, and
$0.4994$ $[0.4592, 0.5395]$ for CREMA-D $\rightarrow$ RAVDESS against
$0.1656$. The oracle values are $0.4555$ $[0.4326, 0.4785]$ and
$0.5680$ $[0.5148, 0.6212]$; \textbf{the oracle is a post-hoc target-selected
benchmark unavailable to the deployable source-only selection protocol}. The
validated-to-oracle difference is the \emph{oracle gap}: $0.1399$
$[0.0231, 0.2568]$ and $0.0687$ $[0.0090, 0.1283]$. It does not identify a
causal effect of source-side selection, because it also includes target-set
maximisation, finite-sample variation and search-space effects.

Figure~\ref{fig:selection} shows the gap per seed, and the forward direction is
where the protocol visibly fails. Forward, \texttt{source\_val} separates
\texttt{none} ($0.7304$) from \texttt{zscore} ($0.7322$) by $0.0018$, while the
same two rungs differ on target by $0.13$. Unlike the source-translation
result, this near-tie is empirical rather than an exact invariance; the
criterion is not symmetric here, merely uninformative at the resolution
available. The consequence is that selection falls back to \texttt{none} on two
of five forward seeds, returning target macro-F1 of $0.2803$ and $0.1508$.
\textbf{The second is below that pair's chance baseline of $0.1665$.} The
protocol uses no target labels and has no source-side indication of this
outcome. The reverse direction selects an aligned rung on all five seeds
(\texttt{mean\_shift} four times, \texttt{zscore} once) and its gap is
correspondingly smaller.

\begin{figure*}[!t]
  \centering
  \includegraphics[width=\textwidth]{validated_vs_oracle.pdf}
  \caption{Validated against oracle target macro-F1 for each seed, with the
    rung the protocol selected annotated on each validated bar. In the forward
    direction the protocol selects the unaligned condition on two of five
    seeds; on one of them the result falls below the chance baseline.}
  \label{fig:selection}
\end{figure*}

This result is limited to selecting the alignment condition. It does not
assess source-side selection for layer aggregation, which
Section~\ref{sec:results-aggregation} addresses separately.

\subsection{The affine comparison as pipeline context}
\label{sec:results-ladder}

The conditions the criterion was selecting among are the affine ladder.
Table~\ref{tab:ladder} gives source-selected target macro-F1 by rung alongside
candidate-row discrepancy summaries, and Table~\ref{tab:primary} gives the
pre-specified comparison family. These are different averaging protocols: the
discrepancy columns include every inner setting, whereas target scores use the
source-selected setting in each cell. Their juxtaposition is descriptive and is
not a matched map-level correlation. The forward summaries include available
transformer cells from its reduced-seed arm, so they are not an equal-family,
identical-candidate-set comparison across seeds.

\begin{table*}[!t]
  \centering
  \caption{The alignment ladder. Within each (pair, seed, backbone, layer aggregation, classifier) cell, the source-validation-best inner setting is selected before target scoring. Target macro-F1 rises once, off \texttt{none}, and then does not track discrepancy: \texttt{mkmmd\_full} reaches the lowest discrepancy in both geometries and is never the best rung.}
  \label{tab:ladder}
  \resizebox{\textwidth}{!}{%
  \begin{tabular}{lrrrrr}
    \toprule
    pair & rung & candidate rows & target macro-F1 & effect (own) & effect (ref) \\
    \midrule
    RAVDESS $\rightarrow$ CREMA-D & \texttt{none} & 153 & 0.2289 [0.2133, 0.2444] & 1432.93 & 36905.40 \\
     & \texttt{zscore} & 153 & 0.3599 [0.3483, 0.3726] & 51.79 & 480.85 \\
     & \texttt{mean\_shift} & 153 & 0.3595 [0.3442, 0.3762] & 104.05 & 19641.52 \\
     & \texttt{coral} & 828 & 0.3648 [0.3480, 0.3831] & 34.84 & 311.15 \\
     & \texttt{mkmmd\_diag} & 558 & 0.3497 [0.3347, 0.3663] & 67.34 & 1057.77 \\
     & \texttt{mkmmd\_full} & 558 & 0.3540 [0.3383, 0.3696] & 12.04 & 23.25 \\
    CREMA-D $\rightarrow$ RAVDESS & \texttt{none} & 135 & 0.2361 [0.1981, 0.2694] & 1045.24 & 23610.38 \\
     & \texttt{zscore} & 135 & 0.4321 [0.4070, 0.4547] & 41.00 & 331.64 \\
     & \texttt{mean\_shift} & 135 & 0.4178 [0.3931, 0.4386] & 76.60 & 8574.50 \\
     & \texttt{coral} & 810 & 0.4193 [0.3968, 0.4394] & 46.55 & 309.77 \\
     & \texttt{mkmmd\_diag} & 540 & 0.4212 [0.3984, 0.4417] & 46.37 & 830.54 \\
     & \texttt{mkmmd\_full} & 540 & 0.3885 [0.3653, 0.4099] & 10.71 & 17.71 \\
    \bottomrule
  \end{tabular}}
  \\[2pt]
  \begin{minipage}{\linewidth}\footnotesize Filter: \texttt{freeze\_tag=grid-freeze-v3}, \texttt{blending=none}. Target intervals are a paired cluster bootstrap over target-test speakers and seeds, 2000 replicates; discrepancy columns are $t$-intervals over seeds. Candidate-row counts differ because CORAL and MK-MMD have larger inner grids; rows are not averaged before selection. Figure~\ref{fig:frames} compares the two geometries.\end{minipage}
\end{table*}


Moving from \texttt{none} to any aligned rung is worth between $+0.1208$ and
$+0.1359$ for RAVDESS $\rightarrow$ CREMA-D and between $+0.1524$ and
$+0.1961$ for CREMA-D $\rightarrow$ RAVDESS. All ten of these contrasts, and
both aggregation contrasts, are in the pre-specified primary family;
\textbf{all 14 survive Holm correction}, with adjusted $p$ at the bootstrap's
resolution floor.

\begin{table*}[!t]
  \centering
  \caption{The pre-specified primary comparisons. All 14 survive Holm correction. The family was fixed in code before the frozen confirmatory grid was run.}
  \label{tab:primary}
  \begin{tabular}{lrrrrr}
    \toprule
    id & comparison & pair & difference in target macro-F1 & Holm $p$ & verdict \\
    \midrule
    L1 & \texttt{zscore - none} & RAVDESS $\rightarrow$ CREMA-D & +0.1310 [+0.1174, +0.1422] & $<$0.0070 & survives \\
    L1 & \texttt{zscore - none} & CREMA-D $\rightarrow$ RAVDESS & +0.1961 [+0.1501, +0.2449] & $<$0.0070 & survives \\
    L2 & \texttt{mean\_shift - none} & RAVDESS $\rightarrow$ CREMA-D & +0.1306 [+0.1159, +0.1439] & $<$0.0070 & survives \\
    L2 & \texttt{mean\_shift - none} & CREMA-D $\rightarrow$ RAVDESS & +0.1817 [+0.1354, +0.2302] & $<$0.0070 & survives \\
    L3 & \texttt{coral - none} & RAVDESS $\rightarrow$ CREMA-D & +0.1359 [+0.1253, +0.1476] & $<$0.0070 & survives \\
    L3 & \texttt{coral - none} & CREMA-D $\rightarrow$ RAVDESS & +0.1833 [+0.1373, +0.2293] & $<$0.0070 & survives \\
    L4 & \texttt{mkmmd\_diag - none} & RAVDESS $\rightarrow$ CREMA-D & +0.1208 [+0.1079, +0.1330] & $<$0.0070 & survives \\
    L4 & \texttt{mkmmd\_diag - none} & CREMA-D $\rightarrow$ RAVDESS & +0.1852 [+0.1406, +0.2317] & $<$0.0070 & survives \\
    L5 & \texttt{mkmmd\_full - none} & RAVDESS $\rightarrow$ CREMA-D & +0.1251 [+0.1149, +0.1351] & $<$0.0070 & survives \\
    L5 & \texttt{mkmmd\_full - none} & CREMA-D $\rightarrow$ RAVDESS & +0.1524 [+0.1057, +0.1991] & $<$0.0070 & survives \\
    A1 & \texttt{layer - last} & RAVDESS $\rightarrow$ CREMA-D & +0.0640 [+0.0541, +0.0748] & $<$0.0070 & survives \\
    A1 & \texttt{layer - last} & CREMA-D $\rightarrow$ RAVDESS & +0.1439 [+0.1140, +0.1746] & $<$0.0070 & survives \\
    A2 & \texttt{weighted - last} & RAVDESS $\rightarrow$ CREMA-D & +0.0597 [+0.0504, +0.0702] & $<$0.0070 & survives \\
    A2 & \texttt{weighted - last} & CREMA-D $\rightarrow$ RAVDESS & +0.1096 [+0.0888, +0.1326] & $<$0.0070 & survives \\
    \bottomrule
  \end{tabular}
  \\[2pt]
  \begin{minipage}{\linewidth}\footnotesize Paired cluster bootstrap over target-test speakers and seeds, 2000 replicates. $p$ values marked $<$ are at the bootstrap's resolution floor of $1/n_{\mathrm{boot}}$ and are not resolvable below it.\end{minipage}
\end{table*}


The among-alignment comparisons are post hoc. No evaluated condition is shown
to outperform \texttt{zscore}. To quantify the remaining uncertainty we use
upper bounds on the four other aligned conditions minus \texttt{zscore},
jointly across both directions; the eight-contrast Bonferroni family uses the
$99.375$th bootstrap percentile, targeting approximate 95\% familywise
coverage. The largest corrected upper bound is $+0.0163$ forward and $+0.0011$
reverse, both determined by CORAL, with point differences of $+0.0049$ and
$-0.0128$. Positive upper bounds are not evidence of an advantage, and these
results do not establish practical equivalence without a predefined
application-specific margin. The conditions can differ in the other direction:
reverse \texttt{mkmmd\_full} minus \texttt{zscore} is $-0.0437$
$[-0.0559, -0.0320]$.

Two limits on reading this table matter for the argument that follows. First,
these scores compare complete pipelines, including the classifier geometries
the maps induce; by Lemma 2 the rungs differ in effective regularisation and
kernel geometry as well as in matched moments, so the table does not establish
that covariance matching has no value or that matching additional moments
causes no benefit. Second, the MK-MMD acceptance check rejects the fitted
candidate in $64.9\%$ of diagonal-map runs and $55.2\%$ of full-map runs, so
the MK-MMD rows describe a finite procedure including its fallback rather than
a converged affine MK-MMD solution. Supplementary Section~S8 records that
procedure in full.

\subsection{The discrepancy--transfer relationship is not measurement-invariant}
\label{sec:results-frames}

Alignment methods are commonly justified by the discrepancy they remove, which
presumes that the relationship between discrepancy and transfer is well defined.
Equation~\ref{eq:kernel-pullback} says it is defined only relative to a
coordinate system. Table~\ref{tab:frames} and Figure~\ref{fig:frames} measure
the consequence on a 13-layer sweep of 2340 runs.

\begin{table*}[!t]
  \centering
  \caption{Measurement-protocol sensitivity. Spearman $\rho$ between a layer's marginal discrepancy and its target macro-F1, across the 13 layers. The adaptive own-geometry and reference-basis calculations give opposite pooled signs.}
  \label{tab:frames}
  \begin{tabular}{lrrr}
    \toprule
    pair & backbone & $\rho$ (adaptive own) & $\rho$ (reference basis) \\
    \midrule
    RAVDESS $\rightarrow$ CREMA-D & \texttt{hubert} & -0.404 & 0.385 \\
     & \texttt{wav2vec2} & -0.372 & 0.224 \\
     & \texttt{wavlm} & -0.514 & 0.230 \\
    CREMA-D $\rightarrow$ RAVDESS & \texttt{hubert} & -0.450 & -0.038 \\
     & \texttt{wav2vec2} & -0.366 & 0.036 \\
     & \texttt{wavlm} & -0.560 & 0.012 \\
    \textbf{pooled} &  & \textbf{-0.444} & \textbf{0.142} \\
    \bottomrule
  \end{tabular}
  \\[2pt]
  \begin{minipage}{\linewidth}\footnotesize Filter: 13-layer sweep, 2340 runs, logreg, 6 rungs, 3 backbones, both directions, 5 seeds. Each correlation is computed across 13 layers within one direction, backbone, rung and seed, then averaged over rungs and seeds. The pooled row averages all 36 cells and 5 seeds. These are descriptive means: shared seeds and features preclude treating cell-and-seed correlations as independent replicates.\end{minipage}
\end{table*}


\begin{figure*}[!t]
  \centering
  \includegraphics[width=\textwidth]{frame_dependence.pdf}
  \caption{Left: Spearman $\rho$ between a layer's marginal discrepancy and its
    target macro-F1 under two stated measurement protocols, for all 36
    (direction $\times$ backbone $\times$ rung) cells. Right: the analytically
    known CORAL map limit. The contrasting trajectories show measurement
    sensitivity; they do not identify one protocol as correct.}
  \label{fig:frames}
\end{figure*}

The mean of the cell-and-seed correlations is $-0.444$ in each rung's adaptive
own geometry and $+0.142$ in the reference-basis calculation. These opposite
signs describe the fixed sweep. Correlations share seeds and feature sets, so an
interval treating all cell-and-seed values as independent is not an inferential
justification for the sign change. Descriptively, 11 of the 36 individual cells
have uncorrected intervals on opposite sides of zero; this count is not used as
an independent significance claim.

The two statistics use the same features, layers and target scores. They differ
in their basis and bandwidth rule. \textbf{The correlation is therefore not
invariant to the stated measurement protocol.} The runs do not identify a
universally correct protocol. They show why a discrepancy--transfer correlation
must state its coordinate basis and bandwidth rule before it can be compared
with another report.

There is a second gap, and it runs the same way. The generalisation bound that
licenses treating discrepancy reduction as progress is stated over a
\emph{classifier-induced} divergence, defined on the symmetric difference of the
hypothesis class (Section~\ref{sec:related-alignment}). A kernel MMD at a fixed
bandwidth is not that quantity: it is a property of the two feature samples
alone and knows nothing about the hypothesis class the learner will search.
Reducing it is therefore not, in itself, reducing the term the theory bounds.

\paragraph{A map limit with known form.} One case fixes the interpretation
independently of any discrepancy measurement. With each covariance regularised
as $C + \varepsilon\,\mathrm{tr}(C)/d \cdot I$, CORAL's map satisfies
\begin{equation}
  M \;\longrightarrow\; \sqrt{\mathrm{tr}(C_t)/\mathrm{tr}(C_s)} \; I
  \label{eq:coral-limit}
\end{equation}
as $\varepsilon$ grows, so the transform collapses to a global scalar rescale
plus a mean shift: \texttt{mean\_shift} with one extra degree of freedom. A
shrinkage probe confirms the approach on real features, with the relative
distance to the nearest scaled identity falling from $0.9254$ to $0.0049$ and
the fitted scalar converging on the predicted limit; supplementary
Section~S5 reports the trajectory. The limiting map is a source-only
transformation, so the shared-scale invariance in
Equation~\ref{eq:median-scale-invariance} does not apply to it. Knowing the
map's form does not fix the measured discrepancy, which is the point.

\paragraph{Recommendation.} We state the consequence as a recommendation rather
than as a characterisation of what other work does or does not report, which we
are not in a position to survey. \textbf{A reported reduction in discrepancy
should state the geometry in which the discrepancy was measured: the kernel
bandwidth and the basis, and how each was fixed.} Where the bandwidth is set by
a data-dependent rule such as the median heuristic, this matters because the
statistic adapts to a global rescaling of both samples
(Equation~\ref{eq:median-scale-invariance}). A source-defined fixed kernel is
one explicit comparison protocol, not a universal remedy; its saturation under
z-score is reported in supplementary Section~S4. An adaptive calculation
answers a different question and should be labelled with its basis and
bandwidth rule.

\subsection{Where source validation is informative}
\label{sec:results-aggregation}

Source-side validation is not uninformative in general, and the same frozen grid
shows it. Comparisons A1 and A2 in Table~\ref{tab:primary} both survive
correction. Replacing final-layer pooling with a fixed mid-stack layer is worth
$+0.0640$ $[+0.0541, +0.0748]$ forward and $+0.1439$ $[+0.1140, +0.1746]$
reverse; replacing it with a learned softmax over all 13 states is worth
$+0.0597$ $[+0.0504, +0.0702]$ and $+0.1096$ $[+0.0888, +0.1326]$. These effects
are larger than any difference \emph{among} the aligned rungs ($\le 0.0437$) and
smaller than the decision to align at all.

Usefulness on one search axis does not imply usefulness on another, and the
contrast is the point: the same criterion that separates layer settings cannot
separate the two translation arms at all. The failure is specific to a decision,
not a property of source validation as such.

\subsection{Robustness of the aligned-versus-unaligned contrast}
\label{sec:results-robustness}

The comparison that the selection result rests on is aligned versus unaligned,
so we report whether it depends on the label harmonisation the two corpora
require. RAVDESS distinguishes \texttt{calm} from \texttt{neutral} and CREMA-D
does not; the main analysis merges them. Two pre-specified controls, one
dropping RAVDESS \texttt{calm} and one additionally excluding \texttt{neutral}
to give a five-class task, both keep every aligned-versus-\texttt{none} contrast
positive, with positive paired 95\% intervals and positive global one-sided
lower bounds under a shared 12-contrast Bonferroni family. Forward improvements
range from $+0.0914$ to $+0.1059$ under the first control and $+0.0612$ to
$+0.0946$ under the second; reverse improvements range from $+0.1121$ to
$+0.1388$ and $+0.1310$ to $+0.1697$. \textbf{The aligned-versus-unaligned
conclusion is therefore not an artefact of the calm-to-neutral mapping.}
Supplementary Section~S1 gives both controls in full.

Three further analyses are reported in the supplement because they qualify the
scope of these results without supporting the central argument. Standard
label-shift correction reduces macro-F1 in 238 of 240 (record, estimator) pairs
despite a small prior divergence, which bounds what a prior-only correction can
do here but does not establish those corrections' assumptions for these corpora
(Section~S2). Class-conditional discrepancy remains nonzero after alignment
under both label mappings; it is a descriptive diagnostic, is not an estimate of
$P(y \mid x)$, and is not offered as an explanation of the similarity among
aligned rungs (Section~S4). A historical blending probe, in which the best
setting of the interpolation parameter is the one that turns blending off,
produced no result that affects the criterion audit (Section~S7). Per-class F1
and the corresponding confusion structure of the source-validated
configuration, which describe how that configuration distributes its errors
rather than what the criterion can identify, are reported in Section~S6.
